% two-column — reading order across two columns, under a title that spans both.
%
% One concern: does the text come out a column at a time, left then right, with the
% paragraph the column break cuts through kept as one paragraph? A reader that takes
% lines across the page interleaves the two columns, and one that breaks at the column
% boundary turns one paragraph into two. The title spans both columns, so it comes
% first and belongs to neither.
%
% The page is small so the text can be short and still put the break inside a
% paragraph: the second paragraph starts in the left column and ends in the right.
% Each paragraph opens by naming its place, so an order defect reads in the diff and
% not only in the score.
\documentclass[twocolumn]{article}
\usepackage{lmodern}
\usepackage[paperwidth=6in,paperheight=3.6in,margin=0.5in]{geometry}
\pagestyle{empty}
% No hyphenation: a word broken at a line end is the hyphenation fixture's concern, not this one's.
\hyphenpenalty=10000 \exhyphenpenalty=10000 \sloppy
\begin{document}
\twocolumn[\section*{Two Columns and a Spanning Title}]

The first paragraph opens the left column. It describes a river that rises in the
mountains, runs north through a wide valley, and reaches the sea after crossing three
lakes and a long stretch of marsh.

The second paragraph begins near the foot of the left column and is long enough to
continue at the top of the right one. It tells how the valley was settled, first by
farmers who drained the lower fields, then by a mill town that grew where the river
narrows, and finally by a railway that followed the old towpath. A reader that splits
this paragraph at the column break has turned one paragraph into two.

The third paragraph closes the right column. It notes that the marsh is now a reserve,
and that the lakes freeze in most winters.
\end{document}
